% theme: speed_direction_force_and_motion
\MajorQuestion{運動の速さと向き／力と運動}
\SetRowSpaceQanda

\Instruction{運動の速さと向きに関する次の問いに答えなさい。}
\QQ{一定時間に物体が移動した距離によって表す量を\NamedBlank{speed}という。}{}
\QQ{\RefBlank{speed}を求めるときに使う、物体が動いた道のりを何というか。}{}
\QQ{最初から最後まで一定の\RefBlank{speed}で進んだとみなして求める\RefBlank{speed}を何というか。}{}
\QQ{ごく短い時間に移動した距離をもとに求める、その時点の\RefBlank{speed}を何というか。}{}
\QQ{物体が進む方向を表す運動の要素を何というか。}{}

\Instruction{運動の記録に関する次の問いに答えなさい。}
\QQ{運動する物体がつけたテープに、一定時間ごとに点を打つ器具を\NamedBlank{record_timer}という。}{}
\QQ{一定時間ごとに強い光を出し、物体の位置を記録する器具を何というか。}{}
\QQ{\RefBlank{record_timer}で記録したテープ上の、となり合う点どうしの間隔を何というか。}{}
\QQ{物体にいつも下向きにはたらく力を\NamedBlank{gravity}という。}{}
\QQ{物体が\RefBlank{gravity}だけを受けて落下する運動を\NamedBlank{free_fall}という。}{}

\Instruction{力がはたらくときの運動に関する次の問いに答えなさい。}
\QQ{運動と同じ向きに一定の力がはたらき、速さが一定の割合で大きくなる運動を何というか。}{}
\QQ{斜面上の物体にはたらく\RefBlank{gravity}のうち、斜面に沿って下向きにはたらく分力を何というか。}{}
\QQ{斜面上の物体にはたらく\RefBlank{gravity}のうち、斜面をおす向きにはたらく分力を何というか。}{}
\QQ{運動の向きと逆向きに一定の力がはたらき、速さが一定の割合で小さくなる運動を何というか。}{}
\QQ{直線上を一定の\RefBlank{speed}で進む運動を\NamedBlank{uniform_motion}という。}{}

\Instruction{慣性と力のおよぼし合いに関する次の問いに答えなさい。}
\QQ{物体がもとの運動の状態を保とうとする性質を\NamedBlank{inertia}という。}{}
\QQ{外から力を加えなければ、物体がもとの状態を保ち続けることを述べた法則を何というか。}{}
\QQ{物体Aから物体Bに力をおよぼすことを\NamedBlank{action}という。}{}
\QQ{\RefBlank{action}と、それに対して相手から同時にはたらき返される力の関係を表す法則を何というか。}{}
\QQ{1つの物体にはたらく2力が打ち消し合い、物体の運動の状態を変えない関係を何というか。}{}
